% !TEX program = xelatex
% ============================================================
%  تقديم مخطوطة — Manuscript Submission Template
%  نموذج تقديم مخطوطة علمية للمجلات مع ترقيم الأسطر (Line Numbers)
%  Compile with: xelatex main.tex (run twice)
% ============================================================
%  Features:
%    - Title block with authors and superscript affiliations
%    - Abstract and keywords
%    - Line numbers via lineno package (\linenumbers)
%    - Sections: Intro, Methods, Results, Discussion, References
%    - Supplementary note placeholder
%    - Fancy headers with running title
% ============================================================

\documentclass[11pt,a4paper]{article}

% ---- Geometry ----
\usepackage[a4paper,margin=2.5cm,top=2.5cm,bottom=2.5cm,headheight=16pt]{geometry}

% ---- Core packages ----
\usepackage{graphicx}
\usepackage{array}
\usepackage{booktabs}
\usepackage{tabularx}
\usepackage{multirow}
\usepackage{enumitem}
\usepackage{float}
\usepackage{caption}
\usepackage{titlesec}
\usepackage{fancyhdr}
\usepackage{setspace}
\usepackage{xcolor}
\usepackage{amsmath}
\usepackage{lineno}

% ---- RTL / Arabic language support (load before hyperref) ----
\usepackage{bidi}
\usepackage[no-math]{fontspec}
\usepackage[quiet]{polyglossia}

% ---- Fonts ----
\setmainfont[Scale=1]{NotoKufiArabic-Regular.ttf}
\newfontfamily\arabicfont[Script=Arabic,Scale=0.95]{NotoKufiArabic-Regular.ttf}
\newfontfamily\englishfont[Scale=0.95]{TeX Gyre Termes}
\setmonofont{NotoKufiArabic-Regular.ttf}
\newfontfamily\arabicfontsf{NotoKufiArabic-Regular.ttf}

% ---- Languages ----
\setdefaultlanguage[calendar=gregorian,hijricorrection=1]{arabic}
\setotherlanguage[variant=british]{english}

% ---- Hyperref ----
\usepackage[breaklinks,hidelinks]{hyperref}
\hypersetup{
  pdftitle={تقديم مخطوطة},
  pdfauthor={اسم المؤلف}
}

% ---- Captions in Arabic ----
\addto\captionsarabic{%
  \renewcommand{\contentsname}{الفهرس}%
  \renewcommand{\figurename}{الشكل}%
  \renewcommand{\tablename}{الجدول}%
  \renewcommand{\refname}{المراجع}%
  \renewcommand{\abstractname}{الملخص}%
}

% ---- Section formatting ----
\titleformat{\section}{\Large\bfseries}{\thesection}{0.7em}{}
\titleformat{\subsection}{\large\bfseries}{\thesubsection}{0.6em}{}
\titlespacing*{\section}{0pt}{1.4ex plus 0.4ex}{0.9ex plus 0.3ex}
\titlespacing*{\subsection}{0pt}{1.0ex plus 0.3ex}{0.7ex plus 0.2ex}

% ---- Headers / footers ----
\pagestyle{fancy}
\fancyhf{}
\fancyhead[R]{\small\itshape عنوان مختصر للمخطوطة}
\fancyhead[L]{\small اسم المؤلف et al.}
\fancyfoot[C]{\small\thepage}
\renewcommand{\headrulewidth}{0.3pt}
\renewcommand{\footrulewidth}{0pt}

% ---- List spacing ----
\setlist[itemize]{topsep=0.3em, itemsep=0.15em, parsep=0pt}
\setlist[enumerate]{topsep=0.3em, itemsep=0.15em, parsep=0pt}

% ---- Table column types ----
\newcolumntype{L}[1]{>{\raggedright\arraybackslash}p{#1}}
\newcolumntype{C}[1]{>{\centering\arraybackslash}p{#1}}

% ---- Line spacing ----
\setstretch{1.4}

% ---- Line numbers for manuscript submission ----
\linenumbers

\begin{document}

% ============================================================
% TITLE BLOCK
% ============================================================
\thispagestyle{fancy}

\begin{center}
{\LARGE\bfseries عنوان المخطوطة (Manuscript Title)}\\[1em]
% TODO: اكتب أسماء المؤلفين هنا مع الانتماءات (Affiliations) كأسس علوية.
{\large
اسم المؤلف الأول\textsuperscript{1}،
اسم المؤلف الثاني\textsuperscript{2}،
اسم المؤلف الثالث\textsuperscript{1,\,3}
}\\[1em]
{\small
\textsuperscript{1} القسم، الجامعة الأولى، المدينة، الدولة\\
\textsuperscript{2} القسم، الجامعة الثانية، المدينة، الدولة\\
\textsuperscript{3} القسم، الجامعة الثالثة، المدينة، الدولة
}\\[0.8em]
% TODO: حدد المؤلف المراسل (Corresponding author).
{\small \textbf{المؤلف المراسل (Corresponding author):} اسم المؤلف الأول}\\
{\small \texttt{email@example.com}}
\end{center}

\vspace{0.5em}

% ============================================================
% ABSTRACT
% ============================================================
\begin{center}
\begin{minipage}{0.9\textwidth}
\noindent\textbf{الملخص (Abstract):}
% TODO: اكتب ملخص المخطوطة هنا (150--250 كلمة). اذكر الخلفية، الهدف، المنهجية، النتائج، والاستنتاج.

\vspace{0.5em}
\noindent\textbf{الكلمات المفتاحية (Keywords):} كلمة مفتاحية 1، كلمة مفتاحية 2، كلمة مفتاحية 3، كلمة مفتاحية 4.
\end{minipage}
\end{center}

\vspace{1em}

% ============================================================
% BODY
% ============================================================

\section{المقدمة (Introduction)}
\label{sec:intro}

% TODO: اكتب مقدمة المخطوطة هنا. وضح المشكلة البحثية وأهميتها.
هذا نص تجريبي للمقدمة. يصف المشكلة البحثية وأهميتها، ويستعرض الأدبيات السابقة ذات الصلة، ويحدد الفجوة المعرفية التي تسعى المخطوطة لسدها.

\section{المنهجية (Methods)}
\label{sec:methods}

% TODO: اكتب منهجية البحث هنا بالتفصيل.
\subsection{تصميم الدراسة (Study Design)}
% TODO: صف تصميم الدراسة.

\subsection{جمع البيانات (Data Collection)}
% TODO: صف طريقة جمع البيانات.

\subsection{التحليل الإحصائي (Statistical Analysis)}
% TODO: صف طرق التحليل الإحصائي.

\section{النتائج (Results)}
\label{sec:results}

% TODO: اكتب النتائج هنا.
هذا نص تجريبي للنتائج. يعرض نتائج التحليل بشكل واضح مع الإشارة إلى الأشكال والجداول.

% ---- Example table ----
\begin{table}[H]
\centering
\caption{ملخص النتائج الرئيسية}
\begin{tabularx}{\textwidth}{L{4cm} X C{3cm}}
\toprule
\textbf{المتغير} & \textbf{الوصف} & \textbf{القيمة} \\
\midrule
% TODO: املأ بيانات الجدول.
المتغير الأول & وصف المتغير & \LR{0.05} \\
المتغير الثاني & وصف المتغير & \LR{0.01} \\
\bottomrule
\end{tabularx}
\end{table}

% ---- Example figure ----
% \begin{figure}[H]
% \centering
% \includegraphics[width=0.8\textwidth]{figure.png}
% \caption{عنوان الشكل}
% \end{figure}

\section{المناقشة (Discussion)}
\label{sec:discussion}

% TODO: اكتب مناقشة النتائج هنا.
هذا نص تجريبي للمناقشة. يفسّر النتائج ويقارنها بالأدبيات السابقة، ويذكر حدود الدراسة.

\section{الاستنتاجات (Conclusions)}
\label{sec:conclusion}

% TODO: اكتب الاستنتاجات هنا.

% ============================================================
% SUPPLEMENTARY NOTE
% ============================================================
\section*{ملاحظات تكميلية (Supplementary Note)}
\addcontentsline{toc}{section}{ملاحظات تكميلية}

% TODO: اكتب الملاحظات التكميلية هنا، أو أشر إلى الملفات التكميلية (Supplementary files).

% ============================================================
% REFERENCES
% ============================================================
\begin{thebibliography}{99}
% TODO: أضف المراجع هنا، أو استخدم ملف BibTeX خارجي.
\bibitem{ref1} اسم المؤلف. (2024). عنوان المرجع. \emph{اسم المجلة}، المجلد، الصفحات.
\bibitem{ref2} اسم المؤلف. (2024). عنوان المرجع. \emph{اسم المجلة}، المجلد، الصفحات.
\bibitem{ref3} اسم المؤلف. (2024). عنوان المرجع. دار النشر، المدينة.
\end{thebibliography}

\end{document}
